\documentclass[11pt]{article}
\usepackage[T1]{fontenc}
\usepackage{lmodern}
\usepackage[margin=1.1in]{geometry}
\usepackage{amsmath, amssymb, amsthm}
\usepackage{booktabs}
\usepackage{xcolor}
\definecolor{linknavy}{RGB}{25, 55, 125}
\usepackage[colorlinks=true, linkcolor=linknavy, citecolor=linknavy, urlcolor=linknavy]{hyperref}
\hypersetup{
  pdftitle={Berge-Fulkerson covers, cycle double covers, flows and exact normality defects of the first counterexamples to the Petersen coloring conjecture},
  pdfauthor={Felipe Santibanez-Leal},
  pdfkeywords={Petersen coloring conjecture, normal edge-coloring, Berge-Fulkerson conjecture, cycle double cover, perfect matching index, oddness, resistance, snarks, SAT, DRAT}
}
\usepackage{orcidlink}
\usepackage{fancyhdr}

\newtheorem{theorem}{Theorem}[section]
\newtheorem{proposition}[theorem]{Proposition}
\newtheorem{lemma}[theorem]{Lemma}
\newtheorem{conjecture}[theorem]{Conjecture}
\theoremstyle{definition}
\newtheorem{definition}[theorem]{Definition}
\theoremstyle{remark}
\newtheorem{remark}[theorem]{Remark}

\newcommand{\PG}{P}
\newcommand{\Gmain}{G_{112}}
\newcommand{\Gsym}{H_{112}}
\newcommand{\Gfty}{G_{52}}
\newcommand{\Gftyb}{G'_{52}}
\newcommand{\Gsix}{G_{68}}
\newcommand{\pd}{\mathrm{pd}}
\newcommand{\ab}{\mathrm{ab}}
\newcommand{\Zfive}{\mathbb{Z}_5}

% Page-1 header block (conventions/manuscript-header-standard.md). The four document
% macros are set by the publishing tool when the version DOI is reserved. At this text
% width the subtype and the two DOIs do not fit on one line of the box, so each has its
% own line.
\newcommand{\doctype}{Preprint}
\newcommand{\docsubtype}{independent certification and consequence audit}
\newcommand{\docversion}{v0.02}
\newcommand{\docdate}{18 September 2026}
\newcommand{\versiondoi}{10.5281/zenodo.22823416}
\newcommand{\conceptdoi}{10.5281/zenodo.22285164}
\newcommand{\docrepo}{github.com/fsantibanezleal/CAOS\_RESEARCH}
\newcommand{\headerblock}{%
\begin{center}
\fbox{\parbox{0.94\textwidth}{\small\raggedright
\textbf{\doctype} (\docsubtype)\\[2pt]
\docversion, \docdate \quad$\cdot$\quad Licence CC BY 4.0\\[2pt]
CAOS Research open-problems programme \quad$\cdot$\quad problem \texttt{petersen-coloring}\\[2pt]
Version DOI \href{https://doi.org/\versiondoi}{\texttt{\versiondoi}}\\[2pt]
Concept DOI (always latest) \href{https://doi.org/\conceptdoi}{\texttt{\conceptdoi}}\\[2pt]
Not peer reviewed. Source code, data and computational records: \href{https://\docrepo}{\texttt{\docrepo}}
}}
\end{center}
\vspace{4pt}}

\title{Berge--Fulkerson Covers, Cycle Double Covers, Flows\\
and Exact Normality Defects of the First\\
Counterexamples to the Petersen Coloring Conjecture}
\author{Felipe Santiba\~nez-Leal\,\orcidlink{0000-0002-0150-3246}\\
\small CAOS Research program, \texttt{github.com/fsantibanezleal/CAOS\_RESEARCH}}
\date{\docdate}

\pagestyle{fancy}
\fancyhf{}
\fancyhead[L]{\small CAOS Research \doctype}
\fancyhead[R]{\small Petersen Coloring Counterexamples: Audit \textperiodcentered\ \docversion}
\fancyfoot[C]{\small \thepage}
\renewcommand{\headrulewidth}{0.4pt}
\fancypagestyle{plain}{\fancyhf{}\fancyfoot[C]{\small \thepage}\renewcommand{\headrulewidth}{0pt}}

\begin{document}
\maketitle
\headerblock

\begin{abstract}
Jaeger's Petersen coloring conjecture (1988), which asserted that every bridgeless cubic graph
admits an edge map onto the Petersen graph sending every vertex star onto a vertex star, was
refuted in August 2026: Putman exhibited two nonisomorphic 112-vertex counterexamples with
SAT-solver certificates, Jooken gave a human-checkable proof, and Goedgebeur, Jooken,
M\'a\v{c}ajov\'a, Mattiolo and Mazzuoccolo found a 52-vertex counterexample together with
infinite families, placing the smallest counterexample between 38 and 52 vertices. The
conjecture had been the strongest known sufficient condition for the Berge--Fulkerson
conjecture and the 5-cycle double cover conjecture. We audit, by exact proof-carrying
computation, what survives on the three retrievable counterexamples. First, independent
encodings sharing no variable scheme with the public ones refute Petersen colorings and normal
5-edge-colorings of all three graphs with proofs checked by \texttt{drat-trim}; five colorable
controls are accepted, and Putman's four public proofs verify under our checker. Second, every
one of the three graphs admits a Berge--Fulkerson cover, a Berge cover by five perfect
matchings, a Fan--Raspaud triple, a 5-cycle double cover and a nowhere-zero 5-flow, all given
as explicit witnesses re-verified from the graph alone; none admits a nowhere-zero 4-flow. Their
perfect matching index is exactly 4, one below the Petersen graph's value. The two 112-vertex
graphs have oddness 4 and resistance 3, whereas the 52-vertex graph has oddness 2 and
resistance 2, every value certified from both sides. All three admit normal and strong normal
6-edge-colorings, so their normal chromatic index is exactly 6. Defining the Petersen defect as
the least number of vertices at which the star condition must fail, we prove that the 52-vertex
graph has defect exactly~2: no single vertex absorbs the obstruction (52 checked refutations),
while every one of its 1{,}326 vertex pairs does (explicit witnesses); the same holds for the two 112-vertex graphs at every one of their 6{,}216 pairs. A parity argument explains the first half in general: for any edge map into the Petersen graph the number of bad vertices is never exactly one, so every counterexample has defect at least~2 and the three known counterexamples are extremal. Finally, we prove that
no counterexample consists solely of copies of the Petersen 4-pole $F$ from which all known
counterexamples are built, and report a counterexample-guided search over compositions of $F$
with free vertices below 52 vertices. All computations are
propositional with checked DRAT certificates or explicit witnesses. The paper decides nothing about the
general conjectures.
\end{abstract}

\section{Introduction}

\subsection{The conjecture and its ladder}

Let $\PG$ denote the Petersen graph. For a cubic graph $G$, a \emph{Petersen coloring} is a map
$\sigma : E(G) \to E(\PG)$ such that for every vertex $v$ of $G$ there is a vertex $w$ of $\PG$
with $\sigma(\partial_G(v)) = \partial_\PG(w)$, where $\partial(\cdot)$ denotes the set of edges
at a vertex; equivalently, the three edges at $v$ map bijectively onto the three edges at $w$.
A proper 5-edge-coloring is \emph{normal} if every edge $uv$ is \emph{poor}
($|c(\partial u) \cup c(\partial v)| = 3$) or \emph{rich} ($= 5$); the \emph{normal chromatic
index} $\chi'_N(G)$ is the least $k$ admitting a normal $k$-edge-coloring.

\begin{conjecture}[Jaeger 1988~\cite{Jaeger1988}]
Every bridgeless cubic graph admits a Petersen coloring.
\end{conjecture}

Jaeger~\cite{Jaeger1985} proved that a cubic graph has a Petersen coloring iff it has a normal
5-edge-coloring (an explicit bijective form is Proposition~2.1 of~\cite{Putman2026}). The
conjecture implies the Berge--Fulkerson conjecture (six perfect matchings covering every edge
exactly twice) and the 5-cycle double cover conjecture~\cite{Jooken2026, MazzMkrt2020}; the
Berge conjecture (perfect matching index at most~5) and the Fan--Raspaud conjecture (three
perfect matchings with empty intersection) follow from Berge--Fulkerson. Every 3-edge-colorable
cubic graph is trivially Petersen colorable, so the whole ladder concerns snarks. Exhaustive
generation had verified the conjecture for every snark on at most 36 vertices~\cite{BGHM2013} and
for every weak snark on 36 vertices~\cite{GMS2019}; an exhaustive check of the weak snarks on 38
vertices shows that any counterexample has at least 40 vertices~\cite{GJMMMU2026}.

\subsection{The disproof}

On 6 August 2026 Putman~\cite{Putman2026} published a simple bridgeless cubic graph $\Gmain$ on
112 vertices and 168 edges with no Petersen coloring and no normal 5-edge-coloring, certified by
CaDiCaL proofs checked with \texttt{drat-trim}, together with a nonisomorphic $D_3$-symmetric
counterexample $\Gsym$; both have girth~5 and edge- and vertex-connectivity~3. The construction
composes copies of the 4-pole $F$ (the Petersen graph minus the endpoints of one edge) with claw
six-poles: $L = 4F + C$ and $\Gmain = 3L + C$. Jooken~\cite{Jooken2026} gave a human-checkable
proof through the Petersen-coloring sets of $F$ and $L$ and a 4-clique obstruction in the line
graph of $\PG$. Goedgebeur, Jooken, M\'a\v{c}ajov\'a, Mattiolo and Mazzuoccolo~\cite{GJMMM2026}
then gave a cyclically 4-edge-connected girth-5 counterexample $\Gfty$ on 52 vertices and infinite
cyclically 4-edge-connected families. The extended account by the same authors and
Ulyanov~\cite{GJMMMU2026} gives a second counterexample $\Gftyb$ on 52 vertices, a purely
theoretical proof for both, counterexamples of every even order at least 60, and the bounds
$40 \le n_{\min} \le 52$ for the order of a smallest counterexample. It leaves open whether every
bridgeless cubic graph has a normal 6-edge-coloring (a conjecture attributed to \v{S}\'amal) and
whether cyclically 5-edge-connected counterexamples exist. A further counterexample $\Gsix$ on 68
vertices, credited there to an announcement on a social network, is available with the other
four from the House of Graphs~\cite{HoG}. For these five graphs~\cite{GJMMMU2026} reports
Berge--Fulkerson covers, perfect matching index at most~4, 5-cycle double covers, strong normal
6-edge-colorings, and proper 5-edge-colorings with exactly two abnormal edges; the corresponding
statements below were obtained independently.

\subsection{What this paper does}

The five counterexamples $\Gmain$, $\Gsym$, $\Gfty$, $\Gftyb$, $\Gsix$ are the first bridgeless
cubic graphs for which the
Petersen route to a Berge--Fulkerson cover or a 5-cycle double cover is closed. We therefore
ask, for each of them, which of the conjectures below the top of the
ladder still hold, and how far each graph is from being colorable. Every answer is exact: a
positive answer is an explicit witness re-verified by a standard-library checker that reads only
the graph, and a negative answer is a DRAT proof checked by \texttt{drat-trim}.
Section~\ref{sec:certify} gives the independent certification (a replication of the published
refutations), Section~\ref{sec:covers} the perfect matching covers, Section~\ref{sec:cdc} the
cycle double covers, flows, oddness and resistance, Section~\ref{sec:normal} the normal
6-edge-colorings, the exact defects, and graphs of arbitrarily large defect with their
consequence for sublinear approximations of the conjecture, Section~\ref{sec:compositions} the
compositions of
copies of the Petersen 4-pole, and Section~\ref{sec:open} the open questions.

\section{Independent certification}\label{sec:certify}

\subsection{The graphs}

We use Putman's normalized digest: the SHA-256 of the compact JSON serialization of the sorted
edge list with $u < v$. The five graphs are
$\Gmain$ (digest \texttt{dc16cc18\ldots{}e8b}), $\Gsym$ (digest \texttt{0f2d8858\ldots{}35c7}),
both transcribed from Putman's CC0 archive; $\Gfty$, transcribed from the appendix
of~\cite{GJMMM2026} with vertices renumbered from $1..52$ to $0..51$ (digest
\texttt{27db5d3b\ldots}); and $\Gftyb$ (digest \texttt{d30a423a\ldots{}1477}) and $\Gsix$ (digest
\texttt{19c7c22a\ldots{}5ab2}), taken from entries 57278 and 57280 of the House of
Graphs~\cite{HoG}. Entries 57237, 57279 and 57244 of that database are isomorphic to $\Gmain$,
$\Gsym$ and $\Gfty$. The full edge lists are in the repository.

\begin{proposition}[machine-verified]\label{prop:structure}
Each of the five graphs is simple, cubic, connected, has edge connectivity~3, girth~5,
and cyclic edge connectivity exactly~4: no set of at most three edges is cycle-separating
(exhaustive search over all edge subsets of size at most three), and an explicit cycle-separating
4-cut exists (edges $\{2,6,9,12\}$ of the sorted edge lists of both 112-vertex graphs,
$\{2,3,8,14\}$ of $\Gfty$, $\{0,11,70,71\}$ of $\Gftyb$ and $\{2,11,90,97\}$ of $\Gsix$). The
automorphism groups of $\Gmain$ and $\Gsix$ are trivial; those of $\Gsym$, $\Gfty$ and $\Gftyb$
have order~6.
\end{proposition}

\subsection{Encodings}

Our Petersen encoding uses only edge-image variables $y_{e,f}$ ($e \in E(G)$, $f \in E(\PG)$), one
image per edge, with the constraint that adjacent edges of $G$ receive distinct adjacent edges of
$\PG$. This suffices because $\PG$ is triangle-free: three pairwise adjacent distinct edges of a
triangle-free graph form a star. Two units break symmetry, using the edge-transitivity of $\PG$
and the transitivity of an edge stabilizer on the four edges adjacent to a fixed edge. Our
normal-5 encoding uses edge-color variables $x_{e,c}$, side-presence variables $p_{e,c}, q_{e,c}$
(``color $c$ appears among the two other edges at $u$'', respectively at $v$), and a rich
indicator $r_e$; richness is $\lnot(p_{e,c} \land q_{e,c})$ for every $c$ and poorness is
$p_{e,c} \leftrightarrow q_{e,c}$ for every $c$. Neither encoding shares a variable scheme with
Putman's vertex-witness and missing-pair encodings.

\begin{theorem}[machine-verified, replication]\label{thm:certify}
None of the five graphs admits a Petersen coloring or a normal 5-edge-coloring. All ten
refutations, and the Petersen refutations of $\Gmain$, $\Gsym$, $\Gfty$ with symmetry breaking
removed, carry DRAT proofs accepted by \texttt{drat-trim} (Table~\ref{tab:certify}). The same encoders accept the
Petersen graph, $K_4$, the prism and the flower snarks $J_5$, $J_7$ with checker-validated
colorings. Putman's four archived proofs are accepted by our \texttt{drat-trim} binary against
his archived formulas.
\end{theorem}

\begin{table}[h]\centering\small
\begin{tabular}{lrrrrr}
\toprule
instance & variables & clauses & solve (s) & proof (bytes) & check (s)\\
\midrule
$\Gmain$ Petersen & 2{,}520 & 73{,}250 & 65.6 & 98{,}426{,}474 & 43.7\\
$\Gmain$ normal-5 & 2{,}688 & 11{,}088 & 1{,}216.8 & 786{,}818{,}036 & 900.5\\
$\Gsym$ Petersen & 2{,}520 & 73{,}250 & 29.5 & 95{,}632{,}373 & 24.5\\
$\Gsym$ normal-5 & 2{,}688 & 11{,}088 & 977.2 & 868{,}156{,}862 & 1{,}293.5\\
$\Gfty$ Petersen & 1{,}170 & 34{,}010 & 1.8 & 4{,}094{,}201 & 1.2\\
$\Gfty$ normal-5 & 1{,}248 & 5{,}148 & 178.8 & 226{,}490{,}830 & 231.8\\
$\Gftyb$ Petersen & 1{,}170 & 34{,}010 & 13.1 & 4{,}038{,}980 & 5.0\\
$\Gftyb$ normal-5 & 1{,}248 & 5{,}148 & 1{,}977.8 & 177{,}769{,}817 & 1{,}476.4\\
$\Gsix$ Petersen & 1{,}530 & 44{,}474 & 18.9 & 9{,}916{,}665 & 10.3\\
$\Gsix$ normal-5 & 1{,}632 & 6{,}732 & 3{,}588.7 & 345{,}057{,}921 & 4{,}492.7\\
\bottomrule
\end{tabular}
\caption{Certified refutations (CaDiCaL 1.7.3 under WSL, \texttt{drat-trim}); every check printed
\texttt{s VERIFIED}. The last four rows were timed while other computations shared the machine.}\label{tab:certify}
\end{table}

\section{Perfect matching covers}\label{sec:covers}

\begin{theorem}[machine-verified]\label{thm:covers}
Each of the five graphs admits a Berge--Fulkerson cover, a Berge cover by five perfect
matchings, three perfect matchings with empty intersection, and a cover by four perfect
matchings. Consequently their perfect matching index equals~4.
\end{theorem}

\begin{proof}
The matchings are explicit (the repository files of computations EXP-002 and EXP-008
list them as edge index sets) and were re-verified by a checker that tests each set to be a perfect
matching and then counts the coverage of every edge. The index is at least~4 because the graphs
are not 3-edge-colorable (Theorem~\ref{thm:certify} and the trivial direction of the conjecture).
\end{proof}

The value~4 is one below the Petersen graph's own perfect matching index~5, which our control
recovers with a verified refutation of the four-matching instance. Thus, on the first
counterexamples, the Berge--Fulkerson, Berge and Fan--Raspaud conjectures all hold, and the
graphs are in this sense better behaved than the Petersen graph.

\section{Cycle double covers, flows, oddness and resistance}\label{sec:cdc}

\begin{theorem}[machine-verified]\label{thm:cdc}
Each of the five graphs admits a 5-cycle double cover (five even subgraphs covering
every edge exactly twice) and a nowhere-zero $\Zfive$-flow, and none admits a nowhere-zero
$\mathbb{Z}_4$-flow. The oddness and resistance are
\[
\mathrm{odd}(\Gmain) = \mathrm{odd}(\Gsym) = 4,\qquad \mathrm{res}(\Gmain) = \mathrm{res}(\Gsym) = 3,
\]
\[
\mathrm{odd} = \mathrm{res} = 2 \quad\text{for } \Gfty,\ \Gftyb \text{ and } \Gsix .
\]
\end{theorem}

\begin{proof}
Double covers and flows are explicit witnesses; the 4-flow refutations carry checked proofs.
Oddness is decided from both sides: for a bound $b$, the formula asks for a perfect matching $M$
and a vertex 2-coloring of the 2-factor $E \setminus M$ with at most $b$ monochromatic edges. An
odd cycle forces at least one monochromatic edge and an even cycle admits none, so the least
satisfiable $b$ is the oddness. Bounds $1, 2, 3$ are refuted with checked proofs for the
112-vertex graphs and bound~$4$ is satisfied by a perfect matching whose 2-factor has exactly four
odd cycles (counted by the checker); for $\Gfty$, $\Gftyb$ and $\Gsix$ bound~$1$ is refuted and
bound~$2$ satisfied.
Resistance uses a deletion indicator per edge and a proper 3-edge-coloring of the remaining
edges, with the same two-sided protocol.
\end{proof}

The oddness~4 of the 112-vertex graphs
is consistent with the theorem that the smallest snarks with oddness at
least~4 and cyclic connectivity~4 have order~44~\cite{GMS2019}. The two 52-vertex graphs and the
68-vertex graph, by contrast, have the least possible oddness of a snark.

\section{Normal 6-edge-colorings and the Petersen defect}\label{sec:normal}

\subsection{Normal chromatic index}

Mazzuoccolo and Mkrtchyan proved that every simple cubic graph admits a normal 7-edge-coloring
and asked whether 6 colors always suffice for bridgeless cubic graphs~\cite{MazzMkrt2020}; the
question is restated as a conjecture in~\cite{GJMMM2026, GJMMMU2026}, where strong normal
6-edge-colorings (every edge rich) of the five graphs are reported.

\begin{theorem}[machine-verified]\label{thm:normal6}
Each of the five graphs admits a normal 6-edge-coloring and a strong normal
6-edge-coloring. Since none admits a normal 5-edge-coloring (Theorem~\ref{thm:certify}), their
normal chromatic index is exactly~6.
\end{theorem}

The colorings are explicit witnesses re-verified by a checker that recomputes, for every edge,
the number of colors on its two end stars. All ten instances solved within two seconds.

\subsection{How far from colorable: the Petersen defect}

Define the \emph{Petersen defect} of a cubic graph $G$ as the least number of vertices at which
the star condition fails, over all maps $E(G) \to E(\PG)$ (equivalently, the least number of
vertices whose star constraints must be dropped for a coloring to exist); a Petersen coloring is
a map of defect~0. Encoding the defect through a cardinality constraint over all vertices proved
hard for the solver (computation EXP-004): the bound-1 instances of $\Gmain$ and $\Gsym$ reached a
30-minute cap without a decision, and the bound-1 instance of $\Gfty$ was refuted in $532.4$
seconds, with a verified proof. We therefore decided the defect by \emph{designated relaxation}:
for each vertex $v$, a formula whose star
constraints are dropped at $v$ only; and for each pair $\{u, v\}$, a formula relaxed at both.

\begin{theorem}[machine-verified]\label{thm:defect52}
The Petersen defect of $\Gfty$ is exactly~2. For every one of its $52$ vertices, the map relaxed
at that vertex alone does not exist (DRAT proofs checked by \texttt{drat-trim}, each in about
two seconds), while every one of the $1{,}326$ vertex pairs admits a map whose star condition
fails exactly at the two relaxed vertices (explicit witnesses; the checker reports exactly those
two bad vertices; the slowest instance took $0.73$ seconds).
\end{theorem}

Thus the obstruction in the smallest known counterexample is not localized: no single vertex,
in particular none of the four vertices outside the six disjoint copies of $F$ (measured in
Section~\ref{sec:compositions}), can absorb it, whereas every pair can. The first half is not
an accident of $\Gfty$:

\begin{theorem}[derived]\label{thm:parity}
Let $G$ be a cubic graph and $\sigma : E(G) \to E(\PG)$ any map. Then the number of vertices of
$G$ at which $\sigma$ fails the star condition is never exactly one. Hence the Petersen defect
of a cubic graph is $0$ or at least $2$, and $\Gfty$ attains the minimum possible defect of a
counterexample.
\end{theorem}

\begin{proof}
For a vertex $v$ let $\chi_v \in \mathbb F_2^{E(\PG)}$ be the indicator of the multiset of the
three labels at $v$ (repeated labels cancel). Every edge of $G$ contributes its label to both
endpoints, so $\sum_{v} \chi_v = 0$. If $u$ satisfies the star condition then $\chi_u$ is the
indicator of a vertex star of $\PG$, an edge cut. Suppose $v$ is the only bad vertex; then
$\chi_v = \sum_{u \ne v} \chi_u$ is a sum of cuts, hence a cut of $\PG$, of odd weight $1$ or
$3$. The Petersen graph is bridgeless, so the weight is $3$ and the three labels at $v$ are
distinct and form a 3-edge cut $\partial(S)$. Since $\PG$ is cyclically 5-edge-connected, one
side of the cut, say $S$, induces a forest with $c \ge 1$ components, so $3 = |\partial(S)| =
3|S| - 2(|S| - c) = |S| + 2c \ge |S| + 2$ and $|S| = 1$: the cut is a star, and $v$ satisfies
the star condition after all, a contradiction.
\end{proof}

The argument gives more: for any map, the sum of $\chi_v$ over the bad vertices lies in the cut
space of $\PG$. For the 112-vertex graphs the theorem already gives defect at least~2; whether
it equals~2 is decided by pair relaxations (computation EXP-006) for the pairs of vertices
outside the twelve disjoint copies of $F$: all $120$ such pairs of
$\Gmain$ and all $120$ of $\Gsym$ are critical (explicit witnesses, at most $3.1$ seconds each),
so both 112-vertex graphs have Petersen defect exactly~2 as well. The full sweeps show
that, as for $\Gfty$, every one of their $6{,}216$ vertex pairs is critical (explicit
witnesses, slowest $20$ seconds; the checker reports exactly the two relaxed vertices as bad in
all $13{,}758$ pair witnesses of the three graphs). We call this property, verified here for the
three graphs and unexplained by the parity argument, \emph{universal 2-criticality}.

\subsection{Abnormal edges and sublinear approximations}\label{sec:abnormal}

For a proper 5-edge-coloring $c$ of a cubic graph $G$, an edge is \emph{abnormal} if it is neither
poor nor rich, and $\ab(G)$ denotes the least number of abnormal edges over all proper
5-edge-colorings of $G$; write $\pd(G)$ for the Petersen defect. Mattiolo, Mazzuoccolo and
Mkrtchyan~\cite{MMM2021} proved that no proper 5-edge-coloring of a cubic graph has exactly one
abnormal edge and asked whether $\ab(G) \le 2$ implies that $G$ has a normal 5-edge-coloring; the
five counterexamples have colorings with exactly two abnormal edges and answer the question in the
negative~\cite{GJMMMU2026}. The two defects are related as follows.

\begin{lemma}[derived]\label{lem:pdab}
$\pd(G) \le \ab(G)$ for every cubic graph $G$.
\end{lemma}

\begin{proof}
View $\PG$ as the Kneser graph on the 2-subsets of $\{1,\dots,5\}$: two subsets are adjacent when
disjoint, and the \emph{color} of the edge $\{A, B\}$ is the element outside $A \cup B$, so that
the three edges at $A$ carry the three colors outside $A$. Let $c$ be a proper 5-edge-coloring with
set of abnormal edges $N$, and for a vertex $v$ let $A_v$ be the complement of $c(\partial(v))$.
For an edge $e = uv$, the \emph{view from $u$} is the edge of $\PG$ at $A_u$ of color $c(e)$. If
$e$ is poor then $A_u = A_v$; if $e$ is rich then $A_u$ and $A_v$ are disjoint and avoid $c(e)$; in
both cases the two views coincide. Let $\sigma(e)$ be the common view if $e$ is normal and the view
from one chosen end if $e$ is abnormal. A vertex all of whose edges carry its own view satisfies
the star condition, so every bad vertex is the end not chosen of an abnormal edge, and there are
at most $|N|$ of them.
\end{proof}

If $c$ has a single abnormal edge $uv$ and the view from $u$ is chosen, then $v$ is bad: the view
from $u$ is an edge at $A_v$ only if $A_u = A_v$ or $A_u \cap A_v = \emptyset$, that is, only if
$uv$ is poor or rich. The map $\sigma$ would then have exactly one bad vertex, so
Theorem~\ref{thm:parity} implies the proposition of~\cite{MMM2021} quoted above. In particular
$\ab(G) \ge \pd(G) \ge 2$ for every cubic graph without a Petersen coloring. All five counterexamples have explicit proper 5-edge-colorings with exactly two abnormal
edges (computations EXP-006 and EXP-008), hence
\[
\pd = \ab = 2 \quad\text{for each of } \Gmain,\ \Gsym,\ \Gfty,\ \Gftyb,\ \Gsix .
\]
For $\Gfty$ the lower bound $\ab \ge 2$ was also confirmed directly: the 42 single-edge
relaxations refuted with checked proofs meet all 14 edge orbits of its automorphism group of
order~6.

Neither defect is bounded. The constructions are those of~\cite{MMM2021}; the argument through the
cut space of $\PG$ gives a stronger conclusion than the one drawn there.

\begin{theorem}[derived]\label{thm:rings}
Let $G_1, \dots, G_t$, $t \ge 2$, be cubic graphs without a Petersen coloring, $e_i = a_ib_i$ an
edge of $G_i$, and let $R$ be obtained from the disjoint union of the graphs $G_i - e_i$ by adding
the edges $a_ib_{i+1}$, indices modulo $t$. Then every map $E(R) \to E(\PG)$ has a bad vertex in
every $V(G_i)$; hence $\ab(R) \ge \pd(R) \ge t$. If every $G_i$ is bridgeless, so is $R$.
\end{theorem}

\begin{proof}
Suppose every vertex of $V(G_i)$ satisfies the star condition under $\sigma$. With the notation of
the proof of Theorem~\ref{thm:parity}, $\sum_{v \in V(G_i)} \chi_v$ is the sum of the indicators of
the labels of the two edges leaving $V(G_i)$, because an edge inside $V(G_i)$ is counted twice.
Each $\chi_v$ is the indicator of a star, so the sum is a cut of $\PG$ with at most two edges, hence
empty, since $\PG$ is 3-edge-connected. The two leaving edges therefore carry the same label $x$,
and giving $e_i$ the label $x$ turns the restriction of $\sigma$ into a Petersen coloring of
$G_i$, a contradiction. If $f$ is an edge of $G_i - e_i$, any two vertices of $V(G_i)$ are joined
in $G_i - f$ by a path, in which $e_i$ can be replaced by the detour around the ring; the added
edges lie on a cycle through all the blocks. So $R$ is bridgeless.
\end{proof}

\begin{theorem}[derived]\label{thm:frames}
Let $K$ be a cubic graph with $t$ vertices and, for each vertex $x$ of $K$, let $G_x$ be a cubic
graph without a Petersen coloring and $v_x$ a vertex of $G_x$. Replace every vertex $x$ of $K$ by
$G_x - v_x$, attaching the three edges of $K$ at $x$ to the three neighbours of $v_x$. Every map
from the edges of the resulting cubic graph to $E(\PG)$ has a bad vertex in every
$V(G_x - v_x)$; hence $\ab \ge \pd \ge t$. If $K$ and all $G_x$ are 3-connected, so is the result.
\end{theorem}

\begin{proof}
If all vertices of $V(G_x - v_x)$ satisfy the star condition, the sum of their vectors $\chi_v$ is
a cut of $\PG$ equal to the sum of the indicators of the labels of the three leaving edges. It has
odd weight, since $|V(G_x)| - 1$ is odd. Weight~1 is impossible in a bridgeless graph, so the
three labels are distinct and form a 3-edge cut of $\PG$, which is a star by the argument in the
proof of Theorem~\ref{thm:parity}. Restoring $v_x$ with these labels on its edges gives a Petersen
coloring of $G_x$. The connectivity statement is the usual substitution along 3-edge cuts.
\end{proof}

In~\cite{MMM2021} it is conjectured that five statements are equivalent: the Petersen coloring
conjecture, and the existence of a sublinear function $f$ such that $\ab(G) \le f(|V(G)|)$ for
all bridgeless, all 2-connected, all 3-connected, respectively all cyclically 4-edge-connected
cubic graphs $G$. It is proved there that the second statement is equivalent to the first, and
that on the last three classes a sublinear bound exists if and only if $\ab$ is bounded by 5, 7
and 9 respectively. The rings of Theorem~\ref{thm:rings} built from $\Gfty$ are 2-connected with
$\ab \ge |V|/52$, and the graphs of Theorem~\ref{thm:frames} built from $\Gfty$ over 3-connected
cubic graphs $K$ are 3-connected with $\ab \ge |V|/51$.

\begin{theorem}[derived]\label{thm:sublinear}
On the class of 2-connected cubic graphs, and on the class of 3-connected cubic graphs, there is
no sublinear function $f$ with $\ab(G) \le f(|V(G)|)$ for all $G$ in the class. More precisely, on
either class such a function exists if and only if every graph of the class has a normal
5-edge-coloring.
\end{theorem}

\begin{proof}
If every graph of the class has a normal 5-edge-coloring, the zero function is a bound. Otherwise
let $G$ be a graph of the class without one. The rings of $t$ copies of $G$ are 2-connected, and
the graphs of Theorem~\ref{thm:frames} over the prisms with $t$ vertices, $t \ge 6$ even, with
every $G_x = G$, are 3-connected when $G$ is; both have $\ab \ge t$ and at most $t\,|V(G)|$
vertices, so $\ab$ is at least linear in the order along these sequences. The graph $\Gfty$ is
3-connected without a normal 5-edge-coloring (Theorem~\ref{thm:certify}).
\end{proof}

Thus the first four of the five statements of~\cite{MMM2021} are false, and their conjectured
equivalence now amounts to the falsity of the fifth: by their Theorem~4, to the existence of a
cyclically 4-edge-connected cubic graph with $\ab \ge 10$. All five known counterexamples are
cyclically 4-edge-connected with $\ab = 2$. A much smaller defect would suffice.

\begin{proposition}[derived]\label{prop:threshold}
If some cyclically 4-edge-connected cubic graph $G^*$ has $\pd(G^*) \ge 3$, then $\pd$ and $\ab$
are unbounded on the class of cyclically 4-edge-connected cubic graphs. Hence on that class either
$\pd \le 2$ everywhere, or no sublinear function bounds $\ab$.
\end{proposition}

\begin{proof}
Let $e_1 = ab$ and $e_2 = cd$ be the end-edges of a path of length three of $G^*$, and let $M$ be
the 4-pole obtained from $G^* - e_1 - e_2$ by attaching a pendant edge at each of $a, b, c, d$.
If $M$ had a map to $E(\PG)$ satisfying the star condition at every vertex, then giving $e_1$ the
label of the pendant edge at $a$ and $e_2$ that of the pendant edge at $c$ would produce a map on
$G^*$ whose only possible bad vertices are $b$ and $d$, so $\pd(G^*) \le 2$. Hence $M$ has no such
map. Join $t$ copies of $G^* - e_1 - e_2$ cyclically as in~\cite{MMM2021}; the result is
cyclically 4-edge-connected~\cite[Proposition~2]{MMM2021}, and every map from its edges to
$E(\PG)$ has a bad vertex in every copy, since a copy without bad vertices would restrict to a map
on $M$. So its defect is at least $t$.
\end{proof}

The cyclic joining of copies of $G - e_1 - e_2$ used
in~\cite{MMM2021} for that class does not force bad vertices by itself: for $\Gfty$ and for
$\Gftyb$, every 4-pole obtained by deleting two independent edges has a Petersen coloring (one
representative per orbit of edge pairs, 482 4-poles for each graph, computation EXP-010), and the boundary
labels of the colorings found are of two kinds only, as the cut space predicts: two crossed equal
pairs, or the four edges adjacent to one edge of $\PG$.

The bounds of Theorems~\ref{thm:rings} and~\ref{thm:frames} are attained on the smallest
instances (computation EXP-009): for the rings $R_t$ of $t$ copies of $\Gfty$ opened at one edge,
$\pd(R_2) = 2$ and $\pd(R_3) = 3$, and for $K_4$ with each vertex replaced by $\Gfty$ minus a
vertex (204 vertices, 3-connected, girth~5) the defect is exactly~4; each witness has exactly one
bad vertex in every copy. For $R_4$ the defect is exactly~4 as well, with a witness found by relaxing one designated vertex in every copy. As a check of Theorem~\ref{thm:rings}, twenty relaxations of
$R_2$ at two vertices of the same copy were refuted with checked proofs.

\section{Compositions of copies of the Petersen 4-pole}\label{sec:compositions}

All known counterexamples contain the 4-pole $F$~\cite{GJMMM2026}. A packing computation
(exhaustive enumeration of the 8-vertex sets with a 4-edge boundary that induce $F$, followed by
an exact maximum disjoint packing) shows that $\Gfty$ contains exactly six disjoint copies of $F$
and four further vertices, and that each 112-vertex graph contains twelve disjoint copies and
sixteen further vertices. This suggests the class $\mathcal C(k, m)$ of simple cubic bridgeless
graphs made of $k$ disjoint copies of $F$ and $m$ free vertices, the $4k + 3m$ semi-edges being
joined by a perfect matching.

\begin{proposition}[derived; machine-witnessed]\label{prop:pure}
Every graph in $\mathcal C(k, 0)$ has a Petersen coloring. Hence no counterexample consists
solely of copies of $F$.
\end{proposition}

\begin{proof}
$F$ admits a Petersen coloring in which its four semi-edges carry the same edge $x$ of $\PG$
(the distance-zero case of Jooken's Lemma~2.1~\cite{Jooken2026}; an explicit such coloring is
recorded in the repository). Color every copy by it, with the same $x$. Every join edge then
carries $x$ at both ends, and every vertex star is a star of a copy of $F$, hence a star of
$\PG$.
\end{proof}

The counterexample-guided search (computation EXP-005: join formula for the matching; structural
check; Petersen-colorability of each candidate; sound blocking clauses learned from every
coloring found, each recorded with its witness) exhausted $\mathcal C(5,0)$ and $\mathcal C(6,0)$
in seconds by exactly such a universal coloring, but did not converge on classes with free
vertices: the 26-vertex control $\mathcal C(3,2)$ stopped at a 15-minute time limit after
$1{,}059$ iterations and $29{,}052$ learned clauses, and $\mathcal C(5,2)$ (42 vertices) stopped at
a two-hour time limit after $1{,}992$ iterations and $80{,}832$ learned clauses, both without
finding a counterexample and without exhausting; $\mathcal C(6,2)$ was not decided. We do not
claim anything about classes the search did not exhaust; a symmetry-broken enumeration is a
natural next step.

\section{Scope and open questions}\label{sec:open}

Nothing here bears on the general Berge--Fulkerson, cycle double cover, 5-flow or normal
6-edge-coloring conjectures: the computations concern five graphs. The questions left open by the
disproof, the order of a smallest counterexample in $[40, 52]$, cyclically 5-edge-connected
counterexamples, and the normal 6-edge-coloring conjecture, remain open. Two questions arise from
Section~\ref{sec:abnormal}. Is there a cyclically 4-edge-connected cubic graph with Petersen defect
at least~3? By Proposition~\ref{prop:threshold} one such graph would show that no sublinear
function bounds $\ab$ on that class, and would complete the proof of the equivalence conjectured
in~\cite{MMM2021}. And is every pair of vertices of every counterexample critical, as it is for
the five known ones?

\section*{Data availability}

{\raggedright
The edge lists, encoders, checkers, witnesses and computation manifests are in the repository
\url{https://github.com/fsantibanezleal/CAOS_RESEARCH}, under
\path{problems/combinatorics/petersen-coloring/}. An identifier EXP-nnn denotes the archived
computation in the folder \path{experiments/EXP-nnn-*} of that directory, with a deterministic
runner and its artifacts with SHA-256 digests. Section~\ref{sec:certify} rests on EXP-001,
Section~\ref{sec:covers} on EXP-002, Section~\ref{sec:cdc} on EXP-003,
Section~\ref{sec:normal} on EXP-004, EXP-006, EXP-009 and EXP-010, and
Section~\ref{sec:compositions} on EXP-005. The values for $\Gftyb$ and $\Gsix$ in
Sections~\ref{sec:covers} to~\ref{sec:normal} come from EXP-008, and their certification in
Section~\ref{sec:certify} from the first part of EXP-007.
\par}

\begin{thebibliography}{99}
\bibitem{Jaeger1988} F.~Jaeger, Nowhere-zero flow problems, in: L.~W. Beineke, R.~J. Wilson
(eds.), Selected Topics in Graph Theory 3, Academic Press, London, 1988, 71--95.
\bibitem{Jaeger1985} F.~Jaeger, On five-edge-colorings of cubic graphs and nowhere-zero flow
problems, Ars Combinatoria 20-B (1985), 229--244.
\bibitem{Putman2026} B.~Putman, A 112-vertex counterexample to the Petersen Coloring Conjecture,
Zenodo v1.1.0 (2026), \href{https://doi.org/10.5281/zenodo.21845291}{doi:10.5281/zenodo.21845291};
arXiv:2608.10012.
\bibitem{Jooken2026} J.~Jooken, A human-checkable proof of the 112-vertex counterexample to the
Petersen coloring conjecture, arXiv:2608.10028v2 (2026).
\bibitem{GJMMM2026} J.~Goedgebeur, J.~Jooken, E.~M\'a\v{c}ajov\'a, D.~Mattiolo, G.~Mazzuoccolo,
A smaller counterexample and infinite family of counterexamples to the Petersen Coloring
Conjecture, Zenodo (2026), \href{https://doi.org/10.5281/zenodo.21933785}{doi:10.5281/zenodo.21933785}.
\bibitem{GJMMMU2026} J.~Goedgebeur, J.~Jooken, E.~M\'a\v{c}ajov\'a, D.~Mattiolo, G.~Mazzuoccolo,
S.~Ulyanov, Disproving the Petersen Coloring Conjecture: theoretical analysis and an infinite
family of counterexamples, arXiv:2608.10028v3 (2026).
\bibitem{MMM2021} D.~Mattiolo, G.~Mazzuoccolo, V.~Mkrtchyan, On sublinear approximations for the
Petersen coloring conjecture, Bull. Inst. Combin. Appl. 92 (2021), 78--90; arXiv:2104.09241.
\bibitem{HoG} K.~Coolsaet, S.~D'hondt, J.~Goedgebeur, House of Graphs 2.0: a database of
interesting graphs and more, Discrete Appl. Math. 325 (2023), 97--107;
\url{https://houseofgraphs.org}.
\bibitem{MazzMkrt2020} G.~Mazzuoccolo, V.~V. Mkrtchyan, Normal edge-colorings of cubic graphs,
J. Graph Theory 94 (2020), 75--91, \href{https://doi.org/10.1002/jgt.22507}{doi:10.1002/jgt.22507}.
\bibitem{BGHM2013} G.~Brinkmann, J.~Goedgebeur, J.~H\"agglund, K.~Markstr\"om, Generation and
properties of snarks, J. Combin. Theory Ser. B 103 (2013), 468--488,
\href{https://doi.org/10.1016/j.jctb.2013.05.001}{doi:10.1016/j.jctb.2013.05.001}.
\bibitem{GMS2019} J.~Goedgebeur, E.~M\'a\v{c}ajov\'a, M.~\v{S}koviera, Smallest snarks with oddness
4 and cyclic connectivity 4 have order 44, Ars Math. Contemp. 16 (2019), 277--298,
\href{https://doi.org/10.26493/1855-3974.1601.e75}{doi:10.26493/1855-3974.1601.e75}.
\bibitem{MMSW2025} Y.~Ma, D.~Mattiolo, E.~Steffen, I.~H. Wolf, Sets of $r$-graphs that color all
$r$-graphs, Combinatorica 45 (2025), Article 16,
\href{https://doi.org/10.1007/s00493-025-00144-4}{doi:10.1007/s00493-025-00144-4}.
\end{thebibliography}

\end{document}
